
\section{Smallest non-zero extensions between two $W$-type modules}
\subsection{Some notions}

{\color{blue} rewrite in the language of Young diagram}
\begin{definition}
An irreducible $\mathbb H_n$-module $\pi$ is said to be calibrated or ladder if the Zelevinsky multisegment 
\[ \left\{ [a_1,b_1],\ldots, [a_n,b_n]\right\}\]
associated to $\pi$ satisfies the
\[    a_1 < a_2 < \ldots < a_n, \quad b_1<b_2< \ldots <b_n .
\]
\end{definition}


\begin{definition}
An irreducible $\mathbb H_n$-module $\omega$ is said to be $W$-lifted if $\omega$ restricted to $S_n$ is also an irreducible $S_n$-representation. 
\end{definition}

We emphasis that the notion of $W$-lifted does not require the module to be unitary, and so the definition is more general one $K$-type representations.


\begin{definition}
Any multisegment $\left\{ [a,b], [a+1, b-1], \ldots, [a+k,b-k]\right\}$ for some $a, b, k \in \mathbb Z$ is said to be uniform.
\end{definition}



\begin{definition}
Let $\pi$ be a ladder $\mathbb H_m$-module with the associated Langlands parameter $\left\{ [a_1, b_1], \ldots, [a_n, b_n]\right\}$.
\begin{enumerate}
\item A ladder $\mathbb H_m$-module $\pi$ is said to be {\it weakly connected} if $a_1 \leq b_2+1, a_2 \leq b_3+1, \ldots , a_{m-1} \leq b_n+1$.
\item A ladder $\mathbb H_m$-module $\pi$ is said to be {\it quais-hook-shaped} if $b_2+1=a_1, b_3+1=a_2, \ldots, b_{n-1}+1=a_n$.
\item The first hook, denoted by $\mathfrak{hk(\pi)}$, of a connected ladder $\mathbb H_m$-module is the collection  \[ \left\{ [a_1,b_1], [a_2,a_2], \ldots, [a_n,a_n] \right\}.\]
\item  The flated hook, denoted by $\Delta_{hk}(\pi)$ of a connected ladder $\mathbb H_m$-module is the segment $[a_n, b_1]$.
\item Let $\pi$ be a weakly connected ladder $\mathbb H_n$-module. We say that $\pi'$ is one-hook-reduced from $\pi$ if the Zelevinsky parameter associated to $\pi'$ is 
\[   \left\{  [a_2+1, b_2], [a_3+1,b_3], \ldots, [a_n+1, b_n] \right\} . 
\]
\item Let $\pi$ be a weakly connected ladder $\mathbb H_n$-module. We say that $\pi$ is {\it connected} if 
\begin{enumerate}
\item $\pi$ is weakly connected; and 
\item we define $\pi_0=\pi$, and recursively define $\pi_i$ to be the hook-reduced from $\pi_{i-1}$, and each $\pi_i$ is weakly connected.
\end{enumerate}
We shall call $\pi_0, \pi_1, \ldots, \pi_k$ to be the hook-reduced sequence from $\pi$.  Here $k$ is the largest possible integer, and in other words, $\pi_k$ is hook-shaped.
\item A connected ladder $\mathbb H_n$-modules is said to be {\it uniform} if for the hook-reduced sequence $\left\{ \pi_0=\pi, \ldots, \pi_k\right\}$ of $\pi$,
\[ \left\{ \Delta_{hk}(\pi_0), \Delta_{hk}(\pi_1), \ldots, \Delta_{hk}(\Delta_k) \right\}
\]
is uniform. We shall call  $\left\{ \Delta_{hk}(\pi_0), \ldots, \Delta_{hk}(\Delta_k) \right\}$ to be the hook multisegment associated to $\pi$.
\end{enumerate}
\end{definition}


\begin{example}
Let $\mathfrak m=\left\{ [0,3], [-3,1]\right\}$ and $\mathfrak m'=\left\{ [0,3], [-3,-2] \right\}$. Then $\mathfrak m$ is connected and $\mathfrak m'$ is not connected.
\end{example}

\begin{definition}
For a multisegment $\mathfrak m=\left\{[a_1,b_1],\ldots, [a_n,b_n]\right\}$ and $c \in \mathbb Z$, let $m(\mathfrak m, c)$ be the number of segments in $\mathfrak m$ such that $a_i\leq c \leq b_i$.
\end{definition}


\begin{lemma} (Steady-no-increasing property) \label{lem steady no inc}
Let $\mathfrak m$ be a uniform multisegment. Then there exists no $c \in \mathbb Z$ such that
\[  m(\mathfrak m, c-1) = m(\mathfrak m, c) < m(\mathfrak m, c+1) .
\]
\end{lemma}




\begin{definition} \label{def uniform ladder add remove}
Let $\pi$ be a uniform connected ladder $\mathbb H_n$-module. Let $\left\{ [a_1,b_1], \ldots, [a_n,b_n]\right\}$ be a multisegment such that $a_1 >a_2>\ldots >a_n$. 
\begin{enumerate}
%\item The middle point of $\pi$ is $a_n+\frac{b_1-a_n+1}{2}$. 
%\item A node $c \in \mathbb Z$ is said to be at the left of $\pi$ if $c <a_n+\frac{b_1-a_n+1}{2}$.
%\item A node $c \in \mathbb Z$ is said to be at the right of $\pi$ if $c>a_n+\frac{b_1-a_n}{2}$. 
\item The left end of a segment $[a_i, b_i]$ is said to be {\it addable} if $a_i-1>a_{i+1}$.
\item The right end of a segment $[a_i, b_i]$ is said to be {\it removable} if $b_i-1>b_{i+1}$. 
\end{enumerate}
\end{definition}



\begin{lemma}
Let $\pi$ be a uniform connected ladder $\mathbb H_n$-module with the associated multisegment $\mathfrak m=\left\{  [a_1,b_1], \ldots, [a_n, b_n]\right\}$. Let $c \in \mathbb Z$. If $a_i=c+1$ for some $i$ and $c \geq b_n$, then the segment $[a_i,b_i]$ is not addable. (not quite precise)
\end{lemma}

\begin{proof} 
Suppose $[c+1, b_i]$ is an addable segment. Note that from the ladder property, any segment $[a_j, b_j]$ other than $[a_i=c+1, b_i]$  satisfying $a_j< c+1\leq b_j$ will then satisfies $a_j\leq c \leq b_j$. On the other hand, there is at most one segment $[a_j,b_j]$ satisfies $a_j \leq c \leq b_j$, but does not satisfy $a_j\leq c+1 \leq b_j$ (as at most one segment takes the form $[a_j,c]$. 

In sum, we must have 
\[  m(\mathfrak m, c) \leq m(\mathfrak m,c+1) .
\]
This contradicts to the uniform property of $\mathfrak m$, as we are assuming $c\geq b_n$.
\end{proof} 



\begin{proposition} \label{prop addable and removable}
We use the notations in Definition \ref{def uniform ladder add remove}. Let $i^*\neq n$ be the largest integer such that the left end of $[a_{i^*}, b_{i^*}]$ is addable. Then there exists a  segment $[a_{j^*}, a_{i^*}-1]$ (for some $j^*$) in $\mathfrak m$ whose right end is removable. 
\end{proposition}


\begin{proof}
In other words, if a segment $[a_i,b_i]$ contains the point $a_{i^*}-1$, then the segment also contains the point $a_{i^*-1}-2$. This implies that 
\[ m(\mathfrak m, a_{i^*}-1) \leq m(\mathfrak m, a_{i^*}-2). \] 

Suppose now there does not exist a removable segment $[a_{j^*}, a_{i^*}-1]$ in $\mathfrak m$. Then we divide into two situations:
\begin{enumerate}
\item[(a)] There exists a segment of the form $[a_{j}, a_{i^*}-1]$. As we have already argued that the segment must be removable, we must also have a segment $[a_{j}, a_{i^*}-2]$. However, with the discussions in the first paragraph, 
\[  m(\mathfrak m, a_{i^*}-1) < m(\mathfrak m, a_{i^*}-2). \] 
\item[(b)] There does not exist a segment of the form $[a_j, a_{i^*}-1]$. Then any segment containing the point $a_{i^*}-1$ also contains the point $a_{i^*}$. However, with the consideration of the segment $[a_{i^*}, b_{i^*}]$. Then 
\[  m(\mathfrak m, a_{i^*}-2)= m(\mathfrak m, a_{i^*}-1)  < m(\mathfrak m, a_{i^*}) .
\]
This violates steady-no-increasing property in Lemma \ref{lem steady no inc}.
\end{enumerate}



\end{proof}

